\chapter{Representación de Riesz y operadores adjuntos}\label{cap:nuevo16}\label{chap:32}
\index{Riesz!representacion en Hilbert@representación en Hilbert}\index{operador!adjunto}

Los desarrollos del capítulo~\ref{cap:nuevo15} describen un vector mediante sus componentes ortogonales. Para estudiar funcionales y operadores necesitamos relacionar esas componentes con el producto interno: una matriz y su transpuesta conjugada describen la misma relación entre vectores vista desde lados opuestos de ese producto. En Hilbert esa relación sigue existiendo sin elegir coordenadas. La herramienta que la hace posible es el teorema de representación de Riesz. La lectura 32 de \cite{valencia} lo utiliza posponiendo su prueba; aquí lo demostraremos antes de aplicarlo. Cerraremos el volumen relacionando los adjuntos con la aproximación de operadores compactos por operadores de rango finito.

\section{Cada funcional continuo es un producto interno}
\label{lectura:32}
\begin{teorema}[Representación de Riesz]\label{ch32:riesz}
Sea $H$ un espacio de Hilbert y $\varphi:H\to\F$ un funcional lineal continuo. Existe un único $v\in H$ tal que
\[
 \varphi(x)=\ip{x}{v}\quad(x\in H),\qquad \norm{\varphi}=\norm v.
\]
\end{teorema}
\begin{proof}
Si $\varphi=0$, tomamos $v=0$. Supongamos $\varphi\ne0$. Su núcleo $M$ es un subespacio cerrado: si $x_n\in M$ converge a $x$, la continuidad da $\varphi(x)=\lim\varphi(x_n)=0$. Elijamos $z$ con $\varphi(z)\ne0$ y sea $r=z-P_Mz$. Entonces $r\in M^{\perp}$ y $\varphi(r)=\varphi(z)\ne0$, por lo que $r\ne0$.

Para cada $x$, el vector $x-\varphi(x)r/\varphi(r)$ pertenece al núcleo. Tomando su producto interno con $r$,
\[
0=\ip{x}{r}-\frac{\varphi(x)}{\varphi(r)}\norm r^{2}.
\]
Así $\varphi(x)=\varphi(r)\ip{x}{r}/\norm r^{2}$. Con nuestra convención de linealidad, el vector representante es
\[
 v=\frac{\overline{\varphi(r)}}{\norm r^{2}}r,
\]
pues el escalar que multiplica a la segunda entrada sale conjugado. Si $v,w$ representan el mismo funcional, $\ip{x}{v-w}=0$ para todo $x$; con $x=v-w$ se obtiene $v=w$.

Cauchy--Schwarz da $|\varphi(x)|\leq\norm x\norm v$, de modo que $\norm{\varphi}\leq\norm v$. Si $v\ne0$, tomando $x=v/\norm v$ obtenemos $|\varphi(x)|=\norm v$, y la desigualdad opuesta. Para $v=0$ las dos normas son cero.
\end{proof}

El representante depende antilinealmente del funcional: el representante de $\alpha\varphi+\beta\psi$ es $\bar\alpha v+\bar\beta w$. Esta observación evita confundir la identificación entre $H$ y $H^*$ con una aplicación compleja lineal.

\section{Formas sesquilineales y operadores}
\begin{definicion}
Para espacios de Hilbert complejos $H_1,H_2$, una forma $\xi:H_1\times H_2\to\C$ es \emph{sesquilineal} si es lineal en la primera variable y antilineal en la segunda. Es \emph{acotada} si existe $C\geq0$ tal que
\[
 |\xi(x,y)|\leq C\norm x\norm y.
\]
\end{definicion}
\index{forma sesquilineal}
Esta condición equivale a continuidad conjunta. En efecto, si existe tal cota,
\[
 |\xi(x',y')-\xi(x,y)|
 \leq C\norm{x'-x}\norm{y'}+C\norm x\norm{y'-y},
\]
que tiende a cero cuando $(x',y')\to(x,y)$. Recíprocamente, la continuidad en $(0,0)$ proporciona $\delta>0$ tal que $|\xi(u,v)|<1$ si $\norm u,\norm v<\delta$. Para $x,y\ne0$ aplicamos esto a $u=\delta x/(2\norm x)$ y $v=\delta y/(2\norm y)$; por sesquilinealidad,
$|\xi(x,y)|\leq4\norm x\norm y/\delta^2$. Si alguno es cero, la desigualdad se cumple porque $\xi$ vale cero.

Si $T\in\Bcal(H_1,H_2)$ y $S\in\Bcal(H_2,H_1)$, las formas $\ip{Tx}{y}$ y $\ip{x}{Sy}$ son acotadas, con cotas $\norm T$ y $\norm S$ por Cauchy--Schwarz y la definición de norma de operador.

\begin{teorema}\label{ch32:formas}
Toda forma sesquilineal acotada $\xi:H_1\times H_2\to\C$ tiene una representación única
\[
 \xi(x,y)=\ip{Tx}{y},\qquad T\in\Bcal(H_1,H_2).
\]
\end{teorema}
\begin{proof}
Para cada $x$ fijo, $\varphi_x(y)=\overline{\xi(x,y)}$ es un funcional \emph{lineal} continuo sobre $H_2$ y $\norm{\varphi_x}\leq C\norm x$. Por Riesz existe un único vector $Tx$ con $\overline{\xi(x,y)}=\ip{y}{Tx}$, o equivalentemente $\xi(x,y)=\ip{Tx}{y}$.

Para $x,z\in H_1$ y $\alpha,\beta\in\C$,
\[
 \ip{T(\alpha x+\beta z)}{y}=\xi(\alpha x+\beta z,y)
 =\ip{\alpha Tx+\beta Tz}{y}\quad\text{para todo }y.
\]
La unicidad del vector representado da $T(\alpha x+\beta z)=\alpha Tx+\beta Tz$. La igualdad de normas de Riesz da $\norm{Tx}=\norm{\varphi_x}\leq C\norm x$, por lo que $T$ es acotado. Si $S$ representara la misma forma, $\ip{(T-S)x}{y}=0$ para todo $y$, y tomando $y=(T-S)x$ se obtiene $Tx=Sx$ para cada $x$.
\end{proof}

\section{Existencia, norma y ejemplos del adjunto}
\begin{teorema}\label{ch32:adjunto}
Para $T\in\Bcal(H_1,H_2)$ existe un único $T^*\in\Bcal(H_2,H_1)$ tal que
\[
 \ip{Tx}{y}_{H_2}=\ip{x}{T^*y}_{H_1}.
\]
Además $\norm{T^*}=\norm T$ y $\norm{T^*T}=\norm T^2$.
\end{teorema}
\begin{proof}
Para cada $y$, el funcional $x\mapsto\ip{Tx}{y}$ es lineal y tiene norma a lo sumo $\norm T\norm y$. Riesz proporciona el vector $T^*y$. Para probar linealidad, notemos que
\[
 \ip{x}{T^*(\alpha y+\beta z)}
 =\bar\alpha\ip{Tx}{y}+\bar\beta\ip{Tx}{z}
 =\ip{x}{\alpha T^*y+\beta T^*z}.
\]
La unicidad implica la igualdad de los segundos argumentos. Riesz también da $\norm{T^*y}\leq\norm T\norm y$, luego $\norm{T^*}\leq\norm T$. La unicidad del adjunto se obtiene de la misma identidad: si $S$ satisface la relación, $\ip{x}{(T^*-S)y}=0$ para todo $x$, y $(T^*-S)y=0$.

Tomar conjugados en la identidad del adjunto da $\ip{T^*y}{x}=\ip{y}{Tx}$; por unicidad, $(T^*)^*=T$. Aplicando la desigualdad de normas recién demostrada a $T^*$ obtenemos $\norm T\leq\norm{T^*}$ y, por tanto, igualdad.

Para cada $x$,
\[
 \norm{Tx}^{2}=\ip{x}{T^*Tx}\leq\norm x\norm{T^*Tx}
 \leq\norm{T^*T}\norm x^{2}.
\]
Tomar el supremo sobre $\norm x=1$ da $\norm T^{2}\leq\norm{T^*T}$. Por otro lado, la desigualdad de composición proporciona
$\norm{T^*T}\leq\norm{T^*}\norm T=\norm T^{2}$. Si $H_1=\{0\}$, todas las normas son cero y no hace falta tomar un supremo sobre una esfera vacía.
\end{proof}

\begin{ejemplo}[Matrices y corrimientos]
En bases ortonormales finitas, si $T$ tiene matriz $(a_{ij})$, entonces $T^*$ tiene matriz $(\bar a_{ji})$: basta expandir $\sum_{i,j}a_{ij}x_j\bar y_i$ y agrupar por $j$ para obtener $\ip{x}{T^*y}$.

En $\ell^{2}(\N)$, sean $R(x_1,x_2,\dots)=(0,x_1,x_2,\dots)$ y $L(x_1,x_2,\dots)=(x_2,x_3,\dots)$. Ambos son acotados y tienen norma uno. Como las series convergen absolutamente,
\[
 \ip{Rx}{y}=\sum_{n=2}^{\infty}x_{n-1}\bar y_n
 =\sum_{n=1}^{\infty}x_n\bar y_{n+1}=\ip{x}{Ly}.
\]
Así $R^*=L$. Se tiene $LR=I$, pero $RLx=x-x_1e_1$. El corrimiento derecho preserva el producto interno y no es suprayectivo; esta distinción será decisiva en el ejercicio~\ref{ej:lectura-32-4}, sobre operadores unitarios.
\end{ejemplo}

\section{Ejercicios de adjunción, normalidad y unitariedad}

\begin{ejercicio}[Lectura 32, ejercicio 1]\label{ej:lectura-32-1}
Probar las reglas (i) $(\alpha S+\beta T)^*=\bar\alpha S^*+\bar\beta T^*$; (ii) $(ST)^*=T^*S^*$; (iii) $T^{**}=T$; (iv) $I^*=I$; (v) $(T^{-1})^*=(T^*)^{-1}$ para $T$ invertible.
\end{ejercicio}


\begin{ejercicio}[Lectura 32, ejercicio 2]\label{ej:lectura-32-2}
\index{operador!normal}
Si $T$ es normal, es decir, $T^*T=TT^*$, probar $\norm{Tx}=\norm{T^*x}$ y que $T-\alpha I$ es normal para todo $\alpha\in\C$.
\end{ejercicio}


\begin{ejercicio}[Lectura 32, ejercicio 3]\label{ej:lectura-32-3}
\index{operador!autoadjunto}\index{operador!positivo}
Probar que $T+T^*$ y $TT^*$ son autoadjuntos, esto es, coinciden con sus adjuntos.
\end{ejercicio}


\begin{ejercicio}[Lectura 32, ejercicio 4: isometrías y unitarios]\label{ej:lectura-32-4}
\index{operador!unitario}\index{isometria@isometría!lineal}
Analizar la afirmación: $\ip{Tx}{Ty}=\ip{x}{y}$ para todo $x,y$ si y solo si $T$ es invertible y $T^{-1}=T^*$.
\end{ejercicio}


\begin{ejercicio}[Lectura 32, ejercicio 5]\label{ej:lectura-32-5}
Probar que la proyección ortogonal $P_W$ sobre un subespacio cerrado es autoadjunta.
\end{ejercicio}


\begin{ejercicio}[Lectura 32, ejercicio 6]\label{ej:lectura-32-6}
Si $T^*T=0$, probar $T=0$. Decidir si la conclusión vale cuando $TT^*=0$.
\end{ejercicio}


\begin{ejercicio}[Lectura 32, ejercicio 7]\label{ej:lectura-32-7}
Demostrar que $\norm{Tv}=\norm T\norm v$ si y solo si $T^*Tv=\norm T^{2}v$.
\end{ejercicio}


\begin{ejercicio}[Lectura 32, ejercicio 8]\label{ej:lectura-32-8}
Probar $\ker T=(T^*(H))^{\perp}$.
\end{ejercicio}


\begin{ejercicio}[Lectura 32, ejercicio 9: polarización de una forma]\label{ej:lectura-32-9}
Sea $\xi$ sesquilineal y $\widehat\xi(v)=\xi(v,v)$. Probar (i)
\[
 \xi(u,v)=\sum_{k=1}^{4}i^k\widehat\xi\left(\frac{u+i^kv}{2}\right),
\]
y (ii) que la forma cuadrática determina la forma sesquilineal.
\end{ejercicio}


\begin{ejercicio}[Lectura 32, ejercicio 10: multiplicación]\label{ej:lectura-32-10}
Sea $h\in C([a,b],\C)$, $a<b$, y $T_hf=hf$ en $C([a,b])$ con producto integral. Hallar (i) el adjunto, (ii) cuándo es normal, (iii) cuándo es autoadjunto y (iv) cuándo es unitario.
\end{ejercicio}


\section{Aproximación de operadores compactos por rango finito}
\index{operador!compacto}\index{operador!rango finito}
Los siguientes resultados enlazan este capítulo con los operadores compactos estudiados anteriormente. La geometría de Hilbert permite una aproximación en norma de operador, más fuerte que aproximar separadamente cada imagen.

\begin{teorema}\label{ch32:compactos}
Si $T:H_1\to H_2$ es compacto, existe una sucesión de operadores acotados de rango finito $T_n$ con $\norm{T_n-T}\to0$. Además, $T$ es compacto si y solo si $T^*$ es compacto.
\end{teorema}
\begin{proof}
Sea $B$ la bola unitaria de $H_1$. La compacidad del cierre de $T(B)$ implica total acotación. Para $\eps>0$, elijamos una red finita $y_1,\dots,y_m$ de radio $\eps$ para $T(B)$, y sea $F=\spanop\{y_1,\dots,y_m\}$. Al ser de dimensión finita, $F$ es cerrado y tiene proyección ortogonal $P_F$. Si $x\in B$, existe $j$ con $\norm{Tx-y_j}<\eps$. Como $P_FTx$ es el mejor aproximante en $F$,
\[
 \norm{Tx-P_FTx}\leq\norm{Tx-y_j}<\eps.
\]
Por tanto $\norm{T-P_FT}\leq\eps$. Elegir $\eps=1/n$ produce $T_n=P_{F_n}T$, de rango finito y acotado. La cota $\norm{P_{F_n}}\leq1$ asegura su acotación también cuando $F_n=\{0\}$.

Si un operador acotado $A:H_1\to H_2$ tiene rango finito, elijamos una base ortonormal $e_1,\dots,e_m$ de ese rango. Para cada $j$, el funcional $x\mapsto\ip{Ax}{e_j}$ tiene, por Riesz, un representante $v_j\in H_1$. Así
\[
 Ax=\sum_{j=1}^{m}\ip{x}{v_j}e_j,
 \qquad A^*y=\sum_{j=1}^{m}\ip{y}{e_j}v_j.
\]
La segunda fórmula se verifica sustituyéndola en la identidad del adjunto y conjugando los coeficientes en la segunda entrada. Su rango está contenido en $\spanop(v_1,\dots,v_m)$, luego $A^*$ también tiene rango finito. Todo operador acotado de rango finito es compacto: la imagen de la bola unitaria está acotada en un espacio de dimensión finita, y su cierre allí es compacto.

La continuidad de la adjunción da $\norm{T_n^*-T^*}=\norm{T_n-T}\to0$. Recordemos por qué un límite en norma de operadores compactos es compacto cuando el espacio de llegada es completo. Si $A_n\to A$ en norma, para una tolerancia $\delta>0$ elegimos $n$ con $\norm{A-A_n}<\delta/3$. Una red finita de radio $\delta/3$ para $A_n(B)$ se convierte, por desigualdad triangular, en una red de radio $2\delta/3$ para $A(B)$. Así $A(B)$ es totalmente acotado. Su cierre es completo y totalmente acotado, por lo que es compacto. Aplicado a $T_n^*$, esto demuestra que $T^*$ es compacto. Finalmente, si $T^*$ es compacto, la dirección recién probada aplicada a él muestra que $T^{**}=T$ es compacto.
\end{proof}

\section{De la geometría de Hilbert al estudio de operadores}
El recorrido desde espacios métricos hasta Hilbert muestra cómo se acumulan las estructuras. La métrica hace posible hablar de aproximación; la completitud asegura límites; la linealidad permite construir y comparar aproximantes; el producto interno identifica el mejor de ellos mediante perpendicularidad. Riesz y los adjuntos convierten esa geometría en un lenguaje para estudiar operadores. Estos resultados serán herramientas de los volúmenes de medida e integración, análisis funcional y análisis armónico.
